\chapter{The Current Transit Landscape: Kepler, K2, and TESS}
\label{chap:kepler_tess}

This chapter sets the observational scene against which \PLATO's
forecast (\Cref{chap:results}) must be read.
\Cref{sec:transit_method} introduces the transit method and the
precision it delivers on small-planet radii.
\Cref{sec:small_planet_demographics} sketches the small-planet
population as \Kepler{}, K2, and \TESS{} have revealed it.  The
subsequent sections then ask what those surveys have actually
delivered in the \HZIED test region --- the short answer being that
they have not delivered anything usable: the current cold-side sample
is so thin that the two-mean Bayes factor is formally undefined.
Along the way I set out the archive sample, its selection biases,
the empty habitable-zone slab, the composition information available,
and the case for \PLATO.

The data, downloader, and exploratory notebook for the
archive-specific parts of this chapter live in
\texttt{HZIED/real\_data/}, deliberately kept separate from the
Bioverse simulator used elsewhere in the thesis so that real-planet
figures cannot leak into the forecast pipeline.


% =======================================================================
\section{Transit method and RV follow-up}
\label{sec:transit_method}

The transit method is the observational technique that has defined
modern exoplanet demographics and underpins every result in this
thesis.  When a planet's orbit is aligned closely enough with our line
of sight, the planet passes in front of its host star once per orbit,
dimming the star's light by a fraction
%
\begin{equation}
    \delta = \left(\frac{R_p}{R_\star}\right)^{\!2},
    \label{eq:transit_depth}
\end{equation}
%
for the duration of the transit
\parencite{Seager2003,Winn2010}.  The transit depth $\delta$ is the
most direct constraint on planet size available across the mass range
accessible to current instruments: the measurement does not require
the planet to be spatially resolved from its host, it scales
favourably with photometric precision, and the same transit can be
observed multiple times over a multi-year mission baseline to build
up signal-to-noise.  The immediate observables --- the orbital period
$P$ (from the interval between successive transits), the depth
$\delta$, and the transit duration $\tau$ --- combined with a set of
stellar parameters ($R_\star$, $M_\star$, and the limb-darkening
coefficients) give the planet radius $R_p$ and the orbital semi-major
axis $a$; the planet mass is inaccessible to transits alone and
requires either radial-velocity follow-up or transit-timing variations
from neighbouring planets in the same system.

Two selection effects intrinsic to the geometry shape every
transit-derived catalogue and both work against the observation of
the \HZIED signature discussed in \Cref{sec:hzied_framework}.  First,
the transit is only visible if the geometric alignment of the orbit
places the planet in front of the star's disk.  For a circular orbit
the probability of such an alignment is approximately $R_\star / a$,
which for a Sun-like host and a habitable-zone orbit is
${\sim}\,0.5\,\%$: long-period planets are geometrically disfavoured
by the same factor by which they are astrophysically interesting.
Second, credible detection requires a signal-to-noise ratio well
above unity accumulated over the mission baseline.  For a stellar
noise floor $\sigma_\mathrm{phot}$ per photometric measurement, the
per-transit SNR scales as
$\delta/(\sigma_\mathrm{phot}/\sqrt{n_\mathrm{cad}})$ where
$n_\mathrm{cad}$ is the number of cadence measurements per transit,
and the total mission SNR scales further as $\sqrt{N_\mathrm{tran}}$
in the number of independent transits.  Small planets on long-period
orbits are thus penalised twice: they present a small $\delta$ (a few
tens of parts per million for an Earth-analog around a solar-type
star) and they contribute only a few transits over any realistic
mission baseline.  The rest of this chapter shows how deeply these
two biases have shaped the current confirmed-planet catalogue, in
which the median instellation sits roughly $400\times$ higher than
the runaway-greenhouse threshold of \Cref{sec:habitable_zone}.

The precision of the derived planet radius depends on $R_\star$ as
much as on $\delta$ itself, and stellar characterisation is therefore
intertwined with planet characterisation.  \PLATO is designed around
this coupling (see \Cref{chap:plato_mission}): its four-year
long-duration observation phase, combined with a target list of
bright asteroseismic stars, propagates to per-planet radius precisions
of a few percent --- the noise floor the \HZIED test requires
\parencite{Schlecker2024}.  The comparable precision on the current
\Kepler catalogue is typically a factor of two to five worse
\parencite{Bryson2021} and dominates the uncertainty budget on
individual small-planet properties as much as any other observational
input.


% =======================================================================
\section{Small-Planet Demographics}
\label{sec:small_planet_demographics}

The starting point for any demographic test on transiting planets is
the underlying occurrence-rate distribution --- what fraction of stars
host what kinds of planets on what orbits.  Two decades of transit
surveys have converged on a picture in which \emph{small} planets,
defined loosely as $R_p \lesssim 4\,\Rearth$, are the numerically
dominant class around dwarf stars of essentially every spectral type.
The Kepler-based occurrence-rate analyses of \textcite{Fressin2013},
\textcite{Kunimoto2020}, and \textcite{Bergsten2022} converge on total
small-planet occurrences of order unity to a few per FGK dwarf at
orbital periods below one hundred days; larger planets, hot Jupiters,
and short-period gas giants are far less common and contribute little
to the numerical count of habitable-zone-relevant systems.

The small-planet population is itself structured.  \textcite{Fulton2017}
first showed that the radius distribution below $4\,\Rearth$ is
bimodal, with two peaks at ${\sim}\,1.3\,\Rearth$ (super-Earths) and
${\sim}\,2.4\,\Rearth$ (sub-Neptunes) separated by a shallow valley
at $R_p \approx 1.6$--$2.0\,\Rearth$: the ``Fulton gap'' or
``radius valley''.  Two competing but not mutually exclusive theories
predict a gap of exactly this shape and location.  The
photoevaporation model of \textcite{OwenWu2017} argues that XUV
irradiation from the young host star strips the H/He envelope from
super-Earth cores, leaving a bare rocky planet below the gap and a
residual envelope-bearing sub-Neptune above it; the core-powered
mass-loss model of \textcite{GinzburgSchlichting2018} makes the same
qualitative prediction from a different physical mechanism, in which
the escape of the primordial envelope is powered by the residual heat
of the cooling planet interior.  Both mechanisms operate on the
${\sim}\,10^{7}$--$10^{8}\,\mathrm{yr}$ timescale of the
pre-main-sequence phase, and their predicted gaps sit within a single
$R_p$ resolution element of what \Kepler observed.  Distinguishing
them requires late-stage measurements that \PLATO approaches but does
not resolve.

For the \HZIED test three quantities from the demographic literature
matter most.  First, the occurrence-rate density of small planets in
the classical habitable zone of solar-like hosts is
$\eta_\oplus \sim 0.2$--$0.6$ per FGK dwarf
\parencite{Bryson2021,Kunimoto2020}, high enough that a well-designed
transit survey around bright FGK stars can accumulate a
habitable-zone-slab sample of tens to hundreds of planets
\parencite{Schlecker2024}.  Second, the mass--radius relationship of
small planets is not deterministic.  \textcite{ChenKipping2017}
showed that the population-averaged relation is probabilistic, with
intrinsic scatter comparable to observational uncertainties, and this
scatter has to be included in any statistical inference on planet
size.  Third, and most importantly for this thesis, the M-dwarf
sub-population appears to be systematically different from the FGK
one.  \textcite{LuquePalle2022} argued from a sample of M-dwarf
transit-plus-RV systems that the small-planet population around M
dwarfs is bimodal in bulk density, with a rocky mode and a water-rich
mode that dominate the demographic in a way not seen around FGK
stars.  This claim was pushed back by \textcite{Rogers2023} and
refined by \textcite{Rogers2024}, but the leading interpretation of
the data still involves a substantial fraction of M-dwarf small
planets carrying non-negligible water inventories --- of order
percent to tens of percent by mass.

The consequence for what follows is that the water inventory of small
planets, treated as a single injection value
$x_{\mathrm{H_2O}}=0.005$ in the original \textcite{Schlecker2024}
framework, should be treated as a \emph{distribution} whose shape may
depend on host spectral type.  \Cref{chap:methodology} adopts four
alternative composition priors motivated by this literature, and
\Cref{chap:results} tests the sensitivity of the \HZIED detection to
each in turn.


% =======================================================================
% Sections 1.5.3–1.5.8 merged into one section for condensed version
\section{Current survey limitations and the case for PLATO}
\label{sec:kt_current_surveys}

The current Kepler/K2/\TESS{} confirmed-planet archive ($N=4\,044$) is
dominated by hot planets: the median instellation is $\sim 400\times$ the
rocky runaway threshold $S_\mathrm{thresh}=280\,\mathrm{W\,m^{-2}}$.
Only eight confirmed planets in the entire archive sit below
$S_\mathrm{thresh}$ at any radius, and only one survives the mildest
\HZIED instellation cut used elsewhere in this thesis, making the
two-mean Bayes factor formally undefined on real data today.  Both
selection effects intrinsic to transit surveys (geometric transit
probability $\propto R_\star/a$ and signal-to-noise
$\propto\sqrt{N_\mathrm{tran}}$) work against the cold side.  The
direct-mass sub-catalogue capable of constraining composition is three
planets deep in the \HZIED window.  The case for \PLATO is therefore
straightforward: its longer baseline, brighter target list, and
FGK-dominated P5 sample are expected to quadruple the cold-side
habitable-zone count and push the cold-to-hot ratio toward unity,
providing the first sample large enough for the SMA-based inference
developed in this thesis.

%% Original sections kept for reference (commented out for condensed doc):
%
% \section{Data source and sample}
\iffalse

Every number in this chapter comes from the NASA Exoplanet Archive
\parencite{Akeson2013}, specifically its `Planetary Systems Composite
Parameters' (\texttt{pscomppars}) table pulled via the archive's TAP
endpoint.  \texttt{pscomppars} carries one row per confirmed planet
with the curators' best parameter set, gap-filling derived quantities
like \texttt{pl\_insol} and \texttt{pl\_dens} where the source
publication left them blank.  Its raw sibling \texttt{ps} keeps every
published parameter set and is a poor choice for a demographic study:
even after asking for one row per planet, only ${\sim}700$
\Kepler{}+K2+\TESS{} entries carry both a radius and an instellation.

Selecting every \Kepler-, K2-, or \TESS-discovered entry in
\texttt{pscomppars} that has both a radius (\texttt{pl\_rade}) and an
instellation (\texttt{pl\_insol}) returns $N_\mathrm{tot} = 4044$
planets, distributed across the three surveys as

\begin{center}
\small
\begin{tabular}{lc}
    \toprule
    Discovery facility & Planets \\
    \midrule
    \Kepler                                        & $2\,766$ \\
    \TESS   & $\phantom{2\,}\,889$ \\
    K2                                             & $\phantom{2\,}\,389$ \\
    \midrule
    Total                                          & $4\,044$ \\
    \bottomrule
\end{tabular}
\end{center}

For the composition discussion in \Cref{sec:kt_composition} the same
query also returns the archive's best-mass estimate \texttt{pl\_bmasse}
with its provenance flag \texttt{pl\_bmassprov} (either
\texttt{'Mass'} for a direct measurement, \texttt{'M-R relationship'}
for a radius-derived value, or \texttt{'Msini'} for a radial-velocity
minimum mass), the bulk density \texttt{pl\_dens}, and the orbital
eccentricity \texttt{pl\_orbeccen}.

The downloader
(\texttt{real\_data/scripts/download\_confirmed\_planets.py}) takes
about three seconds and overwrites the local parquet each time, so
every regeneration of the figures in this chapter reads a fresh
archive snapshot.


% =======================================================================
\section{Sample overview and selection biases}
\label{sec:kt_sample_overview}

Two selection effects shape the joint radius--instellation
distribution of every transit survey, and both work against the
\HZIED test.

\paragraph{Short-period bias.}
Detection efficiency falls with orbital period on two fronts: longer
periods produce fewer transits inside a fixed mission baseline, and
the geometric transit probability itself scales as $R_\star / a
\propto S^{1/2}$.  Together they push the confirmed-planet sample
toward high instellation: more than $80\,\%$ of the \Kepler{}+\TESS{}
sample sits at $S_\mathrm{abs} > 1000\,\mathrm{W\,m^{-2}}$, while the
rocky runaway threshold $S_\mathrm{thresh} = 280\,\mathrm{W\,m^{-2}}$
--- the entire object of interest for this thesis --- sits far down
the cold tail (\Cref{fig:kt_full_scatter}).

\paragraph{The radius gap.}
The small-planet radius distribution is bimodal around
$R_p \approx 1.8\,\Rearth$, the well-studied Fulton gap
\parencite{Fulton2017} usually read as the boundary between rocky
super-Earths and volatile-rich sub-Neptunes.  The bimodality is
clean in the \Kepler subsample; \TESS' brighter, all-sky target list
selects preferentially for larger radii and the K2 sample sits in
between.  The distinction matters because rocky planets are the
physically-motivated \HZIED target: a sub-Neptune already carries a
puffy envelope on either side of $S_\mathrm{thresh}$, so it
contributes to the null-hypothesis plateau rather than to the step
we are trying to detect.

\begin{figure}[H]
    \centering
    \includegraphics[width=\textwidth]{kt_full_sample_rs}
    \caption[Kepler+K2+TESS confirmed planets on the $R_p$--$S_\mathrm{abs}$ plane]{%
        All 4\,044 confirmed \Kepler, K2, and \TESS planets with both a
        measured radius and a computed instellation.  Left: scatter by
        discovery facility (\Kepler blue, K2 green, \TESS red).
        Right: 2D hexbin density with a log colour scale.  Dashed lines
        mark the rocky ($280\,\mathrm{W\,m^{-2}}$) and Hycean
        ($435\,\mathrm{W\,m^{-2}}$) runaway thresholds and Earth's
        instellation ($1361\,\mathrm{W\,m^{-2}}$).}
    \label{fig:kt_full_scatter}
\end{figure}

\Cref{fig:kt_radius_gap} makes both effects visible: the \Kepler
histogram (blue) shows the canonical dip at $R_p \approx
1.6$--$2\,\Rearth$, while the \TESS histogram (red) sits systematically
to its right, and the density plot on the right of
\Cref{fig:kt_full_scatter} shows how sharply the population drops off
below $S_\mathrm{abs}\sim 1000\,\mathrm{W\,m^{-2}}$.

\begin{figure}[H]
    \centering
    \includegraphics[width=\textwidth]{kt_radius_gap}
    \caption[Radius histogram and the Fulton gap]{%
        Radius distribution of the current transit sample.  Left:
        linear-radius zoom on the small-planet regime
        ($R_p\leq 5\,\Rearth$), shaded band marks the Fulton gap
        $1.6$--$2\,\Rearth$.  Right: log-radius histogram covering the
        full range up to gas giants.  The \Kepler sub-sample shows the
        canonical bimodality across the Fulton gap; \TESS is skewed to
        larger radii by its target-selection function.}
    \label{fig:kt_radius_gap}
\end{figure}


% =======================================================================
\section{The habitable-zone slab}
\label{sec:kt_hz_zoom}

Zooming into the \HZIED-relevant window makes the shortage of
cold-side planets concrete.  The median instellation of the whole
$4\,044$-planet sample sits at $S_\mathrm{abs}\approx 1.2\times
10^{5}\,\mathrm{W\,m^{-2}}$, roughly $400$ times the rocky runaway
threshold.  Cutting to $R_p \leq 3\,\Rearth$ and
$S_\mathrm{abs} \in [50,\,1500]\,\mathrm{W\,m^{-2}}$ --- generous
by the standards of the rest of this thesis --- leaves $31$ planets,
all of them on the hot side of $S_\mathrm{thresh}$.  If we go even
wider and drop the radius cut altogether, only $8$ confirmed
\Kepler+K2+\TESS{} planets in the entire archive sit below
$280\,\mathrm{W\,m^{-2}}$; three of those are gas giants
(\Cref{tab:kt_hz_summary}).

\begin{center}
\small
\begin{tabular}{lcc}
    \toprule
    Constraint & $N_\mathrm{cold}$ & $N_\mathrm{hot}$ \\
    \midrule
    Full catalogue & $\phantom{0}\,8$ & $4\,036$ \\
    Any $R_p$, $S \in [50,\,1500]$          & $\phantom{0}\,0$ & $\phantom{0}\,31$ \\
    $R_p \leq 3\,\Rearth$, $S \in [50,\,1500]$ & $\phantom{0}\,0$ & $\phantom{0}\,31$ \\
    $R_p \leq 4\,\Rearth$, $S \in [10,\,3000]$ & $\phantom{0}\,1$ & $\phantom{0}\,108$ \\
    \bottomrule
\end{tabular}
\captionof{table}[Cold/hot split of confirmed planets across cuts]{%
    Cold vs hot counts of \Kepler{}+K2+\TESS{} confirmed planets on
    either side of the rocky runaway threshold
    $S_\mathrm{thresh}=280\,\mathrm{W\,m^{-2}}$, at successively
    relaxed radius and instellation cuts.  The cold-side sample is
    catastrophically thin: only \emph{one} planet (Kepler-1704~b,
    R $\approx$ 1.9\,$\Rearth$) satisfies even the loosest cut used
    elsewhere in this thesis, and none of the standard \HZIED windows
    contain any cold planet at all.}
\label{tab:kt_hz_summary}
\end{center}

\Cref{fig:kt_hz_slab} shows the same slab visually.

\paragraph{The two-mean test cannot be run.}
The two-aggregate Bayes factor of
\Cref{ssec:hypothesis_test_methodology} needs at least a handful of
planets on each side of $S_\mathrm{thresh}$ so that the per-side mean
and its standard error are both defined.  At most one cold planet
survives any of the \HZIED windows we might reasonably choose, so
the test has no data to work on and cannot return a $\ln\Delta Z$.
That is not the same as a weak or inconclusive result: the
observation is that the current sample carries no HZ inner-edge
information at all.  \PLATO changes this by more than an order of
magnitude, as \Cref{sec:kt_platos_role} shows.

\begin{figure}[H]
    \centering
    \includegraphics[width=0.85\textwidth]{kt_hz_slab}
    \caption[HZ-slab scatter for the current transit sample]{%
        Habitable-zone slab ($S\in[50,\,1500]\,\mathrm{W\,m^{-2}}$,
        $R_p\leq 3\,\Rearth$; $N=31$).  Points coloured by discovery
        facility; the cold side ($S \leq 280\,\mathrm{W\,m^{-2}}$) is
        \emph{empty}, so the cold/hot mean-band overlay that this figure
        template supports is suppressed and replaced with the annotation
        that the two-mean test is undefined.  Compare with the P4/P5
        forecast populations in \Cref{fig:kt_plato_overlay}.}
    \label{fig:kt_hz_slab}
\end{figure}


% =======================================================================
\section{Composition hints}
\label{sec:kt_composition}

The archive has no explicit `composition' column, but a few of the
fields it does carry can be turned into composition proxies with more
or less care.  \Cref{tab:kt_composition_coverage} summarises the
options and their coverage.

\begin{center}
\small
\begin{tabular}{lll}
    \toprule
    Column & Coverage & Notes \\
    \midrule
    \texttt{pl\_rade}    & $4\,044\,/\,4\,044$ & direct from the transit \\
    \texttt{pl\_masse}   & $1\,239\,/\,4\,044$ & directly measured (TTV or RV) \\
    \texttt{pl\_bmasse}  & $4\,020\,/\,4\,044$ & `best' mass; may be inferred \\
    \texttt{pl\_bmassprov} & $4\,020\,/\,4\,044$ & provenance flag (see below) \\
    \texttt{pl\_dens}    & $3\,960\,/\,4\,044$ & bulk density, g\,cm$^{-3}$ \\
    \texttt{pl\_orbeccen} & $3\,628\,/\,4\,044$ & orbital eccentricity \\
    \bottomrule
\end{tabular}
\captionof{table}[Composition-relevant columns and their coverage]{%
    Archive columns useful for composition inference and their coverage
    on the Kepler+K2+TESS sample.  The \texttt{pl\_bmassprov} field is
    critical: for $2\,767$ planets the `best' mass is inferred from a
    mass--radius relationship and is therefore circular for
    composition inference; only the $1\,239$ direct-mass planets carry
    true information.}
\label{tab:kt_composition_coverage}
\end{center}

The three composition proxies, in order of coverage but reverse
order of trustworthiness, are the following.

The coarsest is radius alone.  Every planet in the catalogue has
one, and the community convention
\parencite{Fulton2017,ChenKipping2017} sorts them into five bins:
rocky ($R<1.6\,\Rearth$), the Fulton-gap transition
($1.6\,\Rearth\leq R<2\,\Rearth$), sub-Neptune
($2\,\Rearth\leq R<4\,\Rearth$), Neptunian
($4\,\Rearth\leq R<10\,\Rearth$), and gas giant
($R\geq 10\,\Rearth$).  This is complete but crude --- a
$1.5\,\Rearth$ planet is not automatically rocky.

Bulk density adds a bit more information, at least in principle:
iron-enriched at $\rho > 6\,\mathrm{g\,cm^{-3}}$, Earth-like at
$4$--$6$, water-rich or thin-envelope at $2$--$4$,
volatile-dominated below $2$.  The archive has densities for $3\,960$
planets, but $2\,767$ of those are computed from a mass that was
itself inferred from a mass--radius relationship.  Using such a
density to constrain composition is circular; only the direct-mass
subset gives an honest answer.

That direct-mass subset --- planets with
$\texttt{pl\_bmassprov} = \texttt{'Mass'}$, $1\,239$ objects --- is
the gold standard.  With both $M$ and $R$ measured, the mass--radius
diagram is unambiguous, and any planet lying above the pure-water
track of \textcite{Zeng2016,Zeng2019} \emph{must} carry an H/He or
steam envelope simply because water alone cannot make the observed
radius at the observed mass.  Once restricted to the HZ-slab of
\Cref{sec:kt_hz_zoom} the direct-mass sample drops to three planets,
which does not support any statistical statement.  The \HZIED
composition question is genuinely open today, and the ground-based
radial-velocity follow-up around \PLATO's targets is what will
change that.

\Cref{fig:kt_composition_classes} shows the radius classes projected
onto the $R_p$--$S_\mathrm{abs}$ plane; \Cref{fig:kt_composition_mr}
shows the direct-mass mass--radius diagram with the Zeng composition
tracks overlaid.  Both make the same point graphically: the region
where a rocky, cold-side, well-characterised planet would live is
empty.

\begin{figure}[H]
    \centering
    \includegraphics[width=\textwidth]{kt_composition_classes}
    \caption[Radius-based composition class of the current sample]{%
        Left: number of confirmed \Kepler+K2+\TESS planets in each of
        the five radius-derived composition classes.  Right: the same
        classes overlaid on the $R_p$--$S_\mathrm{abs}$ plane, with the
        rocky and Hycean $S_\mathrm{thresh}$ marks retained for
        orientation.  Rocky planets ($R_p<1.6\,\Rearth$) are strongly
        concentrated at very high instellation; the \HZIED-relevant
        cold-side region is empty.}
    \label{fig:kt_composition_classes}
\end{figure}

\begin{figure}[H]
    \centering
    \includegraphics[width=\textwidth]{kt_composition_mr}
    \caption[Mass--radius diagram of the direct-mass subset]{%
        Direct-mass subset ($\texttt{pl\_bmassprov}=\texttt{'Mass'}$;
        $N=1\,239$).  Left: mass--radius diagram with the Zeng+2016
        Earth-like rocky curve, the Zeng+2019 pure-water curve, and a
        rough H/He-envelope reference.  Points are flagged as at or
        below the water curve (blue; may be rocky) versus above it
        (red; must carry an H/He or steam envelope).  Right: the same
        flag projected onto the $R_p$--$S_\mathrm{abs}$ plane.}
    \label{fig:kt_composition_mr}
\end{figure}


% =======================================================================
\section{What PLATO adds}
\label{sec:kt_platos_role}

Putting the current sample and the \PLATO forecast on the same axes
makes the mission's case in one figure.  The confirmed catalogue is
thin exactly where the \HZIED test needs planets --- small radii,
$S_\mathrm{abs}$ of a few hundred W\,m$^{-2}$ --- because those
planets have long periods and produce only a handful of transits over
the missions' baselines.  Longer duration, better photometric
precision, and the FGK-dominated P5 target list are what \PLATO
brings, and \Cref{fig:kt_plato_overlay} shows what that combination
buys us: even the two-year nominal campaign roughly quadruples the
HZ-slab count and, crucially, pushes the cold-to-hot ratio close to
unity.  Both are prerequisites for the SMA-based inference developed
in the rest of this thesis.

\begin{figure}[H]
    \centering
    \includegraphics[width=\textwidth]{kt_plato_overlay}
    \caption[Real vs \PLATO-forecast populations on the $R_p$--$S_\mathrm{abs}$ plane]{%
        Real \Kepler+K2+\TESS confirmed planets (grey) overlaid with
        the simulated \PLATO P4 (M-dwarf, blue) and P5 (FGK, red)
        detectable populations from
        \Cref{sec:hzied_plato}, using the fiducial
        \texttt{agni\_updated}, $x_{\mathrm{H_2O}}=0.005$,
        depth $\geq 60\,\mathrm{ppm}$ configuration.  P4 contributes
        $139$ planets, P5 contributes $266$; both populations sit
        squarely in the currently sparse
        $S_\mathrm{abs}\lesssim 500\,\mathrm{W\,m^{-2}}$ region where
        the \HZIED signature must be measured.}
    \label{fig:kt_plato_overlay}
\end{figure}


% =======================================================================
\section{Chapter summary}
\label{sec:kt_summary}

The transiting-planet catalogue we have today is a hot-planet
catalogue.  Out of $4\,044$ confirmed \Kepler+K2+\TESS{} planets with
both a radius and an instellation, eight sit below
$S_\mathrm{thresh}=280\,\mathrm{W\,m^{-2}}$ at any radius, and only
one survives even the mildest \HZIED cut this thesis uses.  The
simplest diagnostic Bayes factor
(\Cref{ssec:hypothesis_test_methodology}) is undefined on the
sample, and the direct-mass sub-catalogue that could speak to
composition is three planets deep.  These are not gaps to be squeezed
against the priors of a clever inference --- they are why the \HZIED
test has never been tried on real data.  The rest of the thesis
turns to what \PLATO will observe and asks how much of that missing
cold-side sample the step-and-stare campaigns are forecast to
deliver.
\fi
%% End of commented-out original sections

% Figures produced by:
%   Src/figures/kt_full_sample_rs.py
%     → Figures/generated/kt_full_sample_rs.pdf
%   Src/figures/kt_radius_gap.py
%     → Figures/generated/kt_radius_gap.pdf
%   Src/figures/kt_hz_slab.py
%     → Figures/generated/kt_hz_slab.pdf
%   Src/figures/kt_composition_classes.py
%     → Figures/generated/kt_composition_classes.pdf
%   Src/figures/kt_composition_mr.py
%     → Figures/generated/kt_composition_mr.pdf
%   Src/figures/kt_plato_overlay.py
%     → Figures/generated/kt_plato_overlay.pdf
%
% Data source:
%   HZIED/real_data/data/kepler_tess_confirmed.parquet
%   (regenerate with HZIED/real_data/scripts/download_confirmed_planets.py)
